The computational complexity of the proposed algorithm is $O(md_1d_2 + d_2^3)$ at stage 1, while stage 2 requires additional $O(nd_1d_2 + d_1^3)$ computations. This is small compared to the KIV method \citep{Singh2019}, which takes $O(m^3)$ and $O(n^3)$, respectively, and due to the fact that we are using using the inner product deep kernel \eqref{eq:kernel-def}. The computational bottleneck of the KIV method is to compute the inverse of Gram matrix $K=(k(z_i, z_j))_{i,j}$, which is inevitable if we use infinite-dimensional feature map. However, by using a finite-dimensional feature map, our formulation avoids a costly matrix inversion. 

Our method is also substantially faster than DeepIV method \citep{Hartford2017}, although both computational complexities are linear in the data size. In DeepIV, we need to estimate the density of conditional distribution $\expect[X|z]{\vec{\psi}(X)}$\adcomment{$P(X | Z)$}, which is unstable and difficult to parallelize\adcomment{substantiate better this claim, reference?}. On the other hand, the losses $\tilde{\mathcal{L}}_1^{(m)}, \tilde{\mathcal{L}}_2^{(n)}$ in DFIV methods only require us to perform matrix multiplications, which is stable and easy to parallelize.